\chapter{REQUIREMENTS ANALYSIS}

\section{Software Requirements}

\subsection{Python Libraries}

\begin{itemize}
    
    \item \textbf{PyTorch and TensorFlow:} PyTorch is an optimized tensor library useful for deep learning using CPUs and/or GPUs. It provides a industry-standard platform for training and deploying DL models with added tools and libraries. TensorFlow is a similar platform to PyTorch with similar capabilities; in fact, these two exist and are used simultaneously, preferring one over the other for specific tasks. TensorFlow specifically simplifies the ML/DL workflows to run in any environment. For this project which juggles different ML/DL models which need to be trained and deployed, these libraries are extensively used.
    
    \item \textbf{Numpy:} Numpy is marketed as the fundamental Python package/library for scientific computing. It is an interoperable open-source library which is very optimized for multi-dimensional array computations. ML models break-down all elements as large array of numbers and Numpy comes very handy in facilitating all sorts of numerical computations and other mathematical operations on these arrays. It is used in all aspects of image processing in this project.  
    
    \item \textbf{OpenCV:} OpenCV or Open Source Computer Vision Library, as the name suggests, is an open-source Python library that includes various computer-vision algorithms. It defines basic data structures for images, provides modules for image processing, includes functions for reading and writing files in various image formats, and even provides features for working with video and 3D animation. It also offers a hyper-optimized OCR functionality, however this is not applicable for hand-written text. OpenCV working together with Numpy in this project has been used for image processing tasks like thresholding, augmentation, etc. Any time an image is visualized after feeding it to a DL model, OpenCV is most likely used.
    
    \item \textbf{Pandas:} Pandas is an open source data analysis and manipulation tool for Python. It provides labeled data structures (useful for user-defined data too) and statistics-based functions. It is primarily used while working with relational or labeled data (which this project uses a plethora of). 
    
    \item \textbf{Matplotlib and Seaborn} Matplotlib is a Python library for creating static or animated visualizations (graphs, charts). It even allows for interactive graphical representations. All the graphs that showcase performance of trained models have been generated using Matplotlib. Also, in initial testing phase, it was used to generate circuit-related graphs (simulation of say, voltage vs time), first simply static and then later, interactive. Seaborn is a high-level data visualization library based on Matplotlib. It is used to make the visualizations more attractive and informative. It was used in a limited capacity during testing as Matplotlib's features mostly were adequate for this project.
    
    \item \textbf{Scikit-learn:} Scikit-learn comes in during the post-processing stage. It provides tools for predictive data analysis. It provides easy implementation of k-means algorithm (especially useful for this project), among many others. Although used in a rather limited capacity in the project, it proved useful for clustering of junction points. It has also been used for the evaluation of the models during and after training.
    
    \item \textbf{Hugging Face:} Hugging Face provides a large hub of pretrained transformer models for Natural Language Processing (NLP), vision, and more. It is used in the project to integrate LLMs for answering circuit-related queries and working with TrOCR.
    
\end{itemize}

\subsection{YOLOv7}

The YOLOv7 model is used to detect and classify circuit components in hand-drawn circuit images. It outputs bounding boxes, class labels, and confidence scores for each detected component. After comparison with other version of YOLO (v5, v8, v11), YOLOv7 is selected due to its superior detection accuracy, efficient training capabilities, and support for multiple object scales through its enhanced architecture, while also being of an appropriate size in regards to speed and training time.

\subsection{OCR models}

A CRNN-based OCR model trained from scratch, following the DTRB design philosophy for recognizing text is vital for the project as it is the core technology which reads and understands the component values and units of circuit elements. This has been chosen for the project because, foremost, other ready-made OCR models proved to be either too resource-hungry and/or failed to perform well even after some fine-tuning. Additionally, a minimally fine-tuned TrOCR model has also been used. The uses essentially gets to choose between a faster but slightly inferior custom OCR model, or a Microsoft-built robust yet slower TrOCR.

\subsection{KiCad Symbol Library}

KiCad is a free, open-source EDA suite which offers its libraries to be used freely to build any educational project. The project uses the KiCad symbols library, specifically, in order to simply extract the component symbols standardized by KiCad rather than building the icons/representations for components from scratch for use in the frontend schematic.
\subsection{Coding and Training Platforms}

\begin{itemize}
    \item \textbf{Google Colab:} Google Colab is a cloud-based platform that offers a Jupyter notebook environment for developing deep learning systems. For this project, Google Colab provided with Graphics Processing Units (GPUs) for free, which significantly accelerated the training process, enabling experimentation with larger models and datasets. Google Colab has in-built Pytorch and TensorFlow libraries, which made it easy to use them online with minimal setup (barring the occasional, headache-inducing version incompatibility issues). It also allows for Google Drive to be its storage for zipped datasets that need to be used in the notebook.
    
    \item \textbf{Kaggle:} Kaggle is a multi-faceted online platform popular for data science competitions, datasets, code sharing and model training. It was primarily used to find relevant datasets (Circuit Graph Handwritten Diagrams (CGHD), VQA, Handwritten Circuit Diagram (HCD)), code sharing and model training. The dataset of choice for this project was also found in Kaggle. Additionally, it provided a fully hosted environment with popular Python libraries such as Pandas, Numpy, Tensorflow, PyTorch pre-installed to make implementing and training models convenient, especially with its integration with Google Colab. Its free cloud computing service providing 30 hours of GPU use per week proved indispensable in accelerating the training and testing of different models.
    
    \item \textbf{Lightning AI:} Lightning AI is a cloud workspace, similar to Google Colab and Kaggle, which specializes in helping build agents and train models without any tedious setup. Most importantly, through the use of a student account, it allows for the use of extremely powerful GPUs for free (for a reasonable amount of time). Among all the cloud computing services used for the project, Lightning AI proved the most invaluable for implementing and training the models as the GPUs it allotted were super-fast and saved a ton of time during the tedious and time-consuming process.
    
    \item \textbf{Visual Studio Code:} VS Code is the local \gls{ide} everyone on the team settled on, mainly for its auto-completion, syntax highlighting, and copilot integration. Each member extends it further with their own set of extensions tailored to their workflow. All locally executed code runs through VS Code, and its usefulness increased ten-fold once the web development phase began.
\end{itemize}

\subsection{Web Development}

\begin{itemize}
    \item \textbf{Reactjs:} Reactjs serves as an integral technology underlying the dynamic front-end interfaces to be developed for the project. The frontend for this project includes dynamic elements like dragging, erasing, resizing, etc. of the circuits and circuit elements, and Reactjs fulfills these requirements with convenience and speed through the use of its features like virtual Document Object Model (DOM) abstraction, distinction of code into reusable logic and presentation buckets, etc. as well as compatible packages like react-draggable.
    
    \item \textbf{Next.js:} Next.js is a full-stack framework to build web apps. It is built on top of React by adding server-side rendering and static rendering. While Next.js' main selling point is server-side rendering, it also simplifies frontend workflows that are done using React. The client-side interface for this project is built primarily using Next.js.
    
    \item \textbf{FastAPI:} FastAPI is a backend web framework for building APIs based on standard Python type hints. It offers performance on par with NodeJS. It used Starlette (for the web parts), and Pydantic (for the data parts). The backend of this project, especially relating to handling the data, is done using FastAPI and Pydantic. Pydantic is a data validation library for Python that is also powered by type hints, hence making it convenient to be used with FastAPI. Pydantic also checks the data moving between backend and frontend based on predefined structures.
    
    \item \textbf{Uvicorn:} Uvicorn is an ASGI web server implementation for Python. This means it handles HTTP requests asynchronously. The FastAPI standard installation itself includes Uvicorn which is needed to run the application.
    
    \item \textbf{Github:} Github is a proprietary developer platflorm that uses Git to provide distributed version control. This project uses Github for storing and sharing the project code and resources among project members.
\end{itemize}

\section{Feasibility Study}

\subsection{Technical Feasibility}

Technical Feasibility is assessing if a project idea can be built with the team's current (or selected) technology and resources. It also assesses whether or not the project is within the skill-level of the project team. Foremost, the technology used for the project has been selected after consideration of their features and support. Python libraries such as TensorFlow, Numpy, OpenCV, Matplotlib, etc., React and Netxt.js for frontend, FastAPI with Pydantic and Uvicorn for backend, are all mature, ever-developing technologies and can handle the requirements for this project. They have been selected with the expectation that these technologies will be supported for a considerable time in the future. They also have excellent docs and community support behind them, useful for handling the inevitable issues that arise during development and testing. The project members are decent developers with considerable experience in research and Python-based developement. Overall, the project is technically feasible.

\subsection{Economic Feasibility}

The total expenses for this project have come out to be NPR 41,160. The biggest part of this is the subscription for Google Colab Pro and Lightning AI credits, both useful in training ML/DL models (NPR 24,000). Lightning AI provides token-based scaling and also a generous free tier which is why it was used more than Google Colab Pro. Plus, a lot of research has gone into this project, and the scientific papers required for this were made available through an IEEE Xplore Membership amounting to NPR 2,160. Miscellaneous costs such as report printing and presentation costs as well as API costs, inevitable to any software-based research project amount to 15,000. The total cost is considered justified for a project of this scale and caliber, and hence is economially feasible. Gantt Chart in the Appendix expands on this.

\subsection{Time Feasibility}

The project has taken almost a year to complete by a team of four members. This was expected as it has a lot of moving parts and modules. Before any development, around a month was dedicated to preliminary research and literature review. The scope of this project changed in subtle ways during this, depending on the shortcomings of the existing research and products. This research, however, did not stop once development started, rather, it took a background role throughout the project, helping shape the project.

Phased development approach was followed. Phase A covered the circuit recognition (YOLO), text extraction (custom-OCR and TrOCR), and netlist generation. Phase A was developed and tested by September of this year. A buffer period was planned between Phases A and B for adjustments.

Phase B broadened the feature list. Integration of all elements of Phase A was done into a barebones web application. CircuitJS was roped in for handling the editing of schema and simulation needs of the project. DeepSeek v3.1 was utilized for the LLM-based query feature. Gradual, iterative developments in the UI/UX and subsequent testing continued in the background. This phased approach kept the design intact while layering on the more advanced features over time.

As of yet, all of the promised features have been delivered with only minor kinks to be sorted out. The project is considered feasible from a time perspective.

